% 第七章：用户确认稿缓存，2026-10-16。章节片段，最终五页稿另行归并排版。
\section{总结}
本作品面向几何内核中以样条形式存储的解析几何恢复需求，提出一套以\textbf{目标母面 NURBS 参数域中的线性 PCurve}为边界约束的转换方案。方案覆盖平面、圆柱面、圆锥面和球面，并将规则圆环面作为扩展目标，将母面类型识别、几何参数恢复与有限边界重建纳入统一流程，为后续采用解析表示进行存储和几何运算提供基础。

本方案的主要贡献在于将\textbf{结构驱动的求解与可追溯的几何验收相结合}。等参专用分支优先利用明确的几何结构，一般兜底通过多候选搜索和必要的参数优化扩展处理范围；齐次 Bernstein 复合将线性 PCurve 转化为可计算的三维有理边界；正比例切向残差与分层区间求界则为连续边界验收提供依据。针对实际内核输出曲线重新构造验收数据，使参数求解、边界构造和最终几何表示之间形成完整的验证链条。

鲁棒性体现在对退化情形和失败状态的明确处理。方案区分球面极点处的参数退化、圆锥顶点处的几何奇异，以及正则圆环面上的裁剪边坍缩，依据相应条件保留位置验收并处理切向适用范围。同时，两条计算路径共享时间预算，对无法认证或预算耗尽的输入返回明确失败状态，避免将未完成的判定作为有效转换结果。

定向测试中，三个代表案例均通过转换与独立输出审计，各例三次测试的最大整体耗时均小于 $10\,\mathrm{ms}$，展示了方案在有限面片恢复和退化边界处理方面的实际能力。作品形成了具有明确适用条件、独立验收机制和可复现测试依据的工程实现，为 CAD 几何内核中的解析几何恢复提供了可集成的算法基础。
